% Section 8.1: From Resolved Sources to the Cosmic Web
% Sources: 8.1_introduction.md, 8.1.1_thesis_so_far.md, 8.1.2_dark_matter_cosmic_web.md

\section{From Resolved Sources to the Cosmic Web}
\label{sec:8.1}

Dark matter leaves no preferred address.
It annihilates wherever it finds itself dense enough, which means in every gravitationally bound structure from the Milky Way's central parsecs to the faintest galaxy-cluster subhalos a billion light-years away.
The preceding chapters have probed this emission one scale at a time: the Galactic Center at sub-degree resolution, individual subhalos as point sources, and the statistical aggregate of the sub-threshold population.
Each approach identified a genuine signal channel and, in doing so, also identified what it could not access.
This section steps back to map that landscape explicitly, and then to introduce the complementary strategy that forms the core of Part~IV: exploiting the spatial cross-correlations between the unresolved gamma-ray sky and the large-scale structure of the universe as traced by galaxy catalogs.

Section~\ref{sec:8.1.1} reviews what Parts~II and III established and where the limiting systematics lie. Section~\ref{sec:8.1.2} then develops the physical argument for why the dark matter--large-scale-structure cross-correlation provides orthogonal sensitivity: it selects a low-redshift emission window where dark matter dominates astrophysical backgrounds, and it suppresses noise from source populations that do not cluster with local galaxies.

\subsection{The Thesis So Far: Scales and Strategies}
\label{sec:8.1.1}

The work collected in Parts~II and III progressively refined both the targets and the statistical methods of dark matter searches in gamma rays.
Each stage addressed a specific astrophysical scale and, in doing so, identified the limitations that motivated the next approach.
A brief summary of what each analysis achieved --- and where it fell short --- sets the stage for the strategy taken in this final part.

The Galactic Center excess (GCE), analyzed in Chapter~\ref{ch:04}, represents the strongest expected signal from dark matter annihilation in the Milky Way: an extended, approximately spherically symmetric excess of GeV photons consistent with annihilation by a $\sim35$--50 GeV WIMP into $b\bar{b}$ with a cross-section near the thermal relic value~\cite{Daylan:2014rsa}.
The signal morphology --- scaling roughly as the square of a generalized NFW density profile with inner slope $\gamma \approx 1.1$--1.3, extending to at least $10^\circ$ from the Galactic Center --- matches dark matter expectations well.
Yet a decade of refined analyses has not resolved its origin.
The central problem is that an unresolved population of millisecond pulsars (MSPs) in the Galactic bulge would produce spectral and morphological signatures nearly indistinguishable from those of a smooth dark matter signal~\cite{Bartels:2015aea}. Early statistical arguments using Non-Poissonian Template Fitting (NPTF) appeared to favor a point-source origin for the excess; subsequent investigations showed those results are susceptible to systematic effects --- unmodeled north-south asymmetries in the data can cause a smooth dark matter signal to be spuriously attributed to point sources~\cite{Leane:2019xiy}. Deep learning approaches have partially circumvented these biases, but they encounter the same fundamental degeneracy: in the limit where individual sources contribute much less than one photon each to the dataset, Poissonian (dark matter-like) and non-Poissonian (point-source-like) emission become statistically indistinguishable~\cite{Chang:2019ars}.
The GCE remains unresolved, and breaking the degeneracy will require either the direct radio detection of a bulge MSP population by next-generation telescopes or a corroborating dark matter signal from an independent target.

The search for individual dark matter subhalos among \textit{Fermi}-LAT unassociated sources, described in Chapter~\ref{ch:05}, approached the problem from a very different angle: rather than modeling a known extended excess, it asked whether any of the 1282 unassociated 4FGL-DR4 sources at high Galactic latitude ($|b| > 10^\circ$) belong to the dark matter subhalo population predicted by $\Lambda$CDM. We constructed the first generative mixture model of unassociated sources, representing them as a combination of astrophysical components (active galactic nuclei, pulsars, and other Galactic classes) plus a contribution from dark matter subhalos whose J-factor distribution was drawn from N-body simulations~\cite{Amerio:2025fhz}.
The key methodological advance was the explicit treatment of dataset shift: the distributions of features for associated and unassociated sources differ both in class prevalence (prior shift) and in feature space (covariate shift), and previous classify-and-count methods ignored both effects.
Our maximum likelihood analysis found no significant excess beyond the astrophysical components.
For the $b\bar{b}$ annihilation channel, 95\% confidence upper limits on $\langle\sigma v\rangle$ were derived for dark matter masses from 10 GeV to 1 TeV, competitive with previous searches~\cite{Coronado-Blazquez:2021avx} and consistent with the null result at current \textit{Fermi}-LAT sensitivity.

Chapters~\ref{ch:06} and \ref{ch:07} shifted scope from individual source identification to the aggregate statistical properties of the gamma-ray source population.
Rather than searching for dark matter signatures in specific sources, these analyses characterized the entire sky below the detection threshold.
The source-count distribution $dN/dS$ was recovered using simulation-based inference with a deep convolutional neural network trained on $\sim10^6$ synthetic \textit{Fermi}-LAT maps~\cite{Amerio:2023uet}, extending the measurement of the flux distribution to roughly 50 times below the \textit{Fermi}-LAT detection limit --- a factor-of-ten improvement over previous 1pPDF methods~\cite{Zechlin:2015wdz}. Building on this, the probabilistic catalog of Chapter~\ref{ch:07} combined the recovered $dN/dS$ with a frequentist simulation-comparison framework to assign spatially resolved quality factors to sub-threshold directions in the sky, identifying $\sim50\%$ more candidate source positions than the standard 4FGL catalog~\cite{Amerio:2023rjn}.
Together, these analyses characterized in detail the astrophysical landscape in which any dark matter contribution must reside.
They did not target dark matter specifically, but the population information they provide --- in particular the globally consistent significance scale of the probabilistic catalog --- feeds directly into any cross-correlation study that requires a clean characterization of the unresolved sky.

The conceptual thread connecting Parts~II and III is one of increasing scale and decreasing signal strength.
The GCE offers the highest expected dark matter surface brightness but is contaminated by Galactic systematics and thermal-emission degeneracies.
Subhalo searches target individual objects at predictable fluxes but face the sparse-signal problem: the expected number of detectable dark matter subhalos at current sensitivity is at most of order a few dozen, and none was observed.
Population statistics describe the aggregate unresolved sky but do not isolate a dark matter component.
The signal we have not yet exploited is one that is distributed across all angular scales: the spatial correlation pattern that the hierarchical distribution of dark matter halos --- from the nearest cluster down to the smallest subhalo --- imprints on the gamma-ray sky at cosmological distances.

\subsection{Dark Matter in the Cosmic Web}
\label{sec:8.1.2}

If dark matter is a particle that annihilates, it does so everywhere it is dense enough for pairs to meet --- not just in the Galactic Center or in individual subhalos, but inside every gravitationally bound structure across the observable universe.
Understanding how much gamma-ray emission this process produces at cosmological distances, and how that emission is distributed across the sky, requires first understanding how dark matter structures are arranged in space.

In the standard $\Lambda$CDM cosmological framework, structures form hierarchically from the bottom up.
Small-amplitude density perturbations in the early universe collapse first into low-mass dark matter halos; these merge and accrete material to build progressively more massive structures over cosmic time.
The result is a cosmic hierarchy: galaxy clusters at the top, hosting individual galaxy halos, which in turn host populations of subhalos down to the minimum halo mass set by the free-streaming length of the dark matter particle --- roughly Earth-mass ($\sim 10^{-6} \, M_\odot$) for a typical $\sim100$ GeV WIMP~\cite{Fornasa:2015qua}.
This hierarchical web of halos and subhalos, distributed across the universe, is what we call the cosmic web of dark matter.

The gamma-ray flux from dark matter annihilation is proportional to the square of the local dark matter density.
Dense, compact structures are therefore the dominant contributors to the total signal, and the cumulative emission from all unresolved halos and subhalos across all redshifts produces a diffuse, approximately isotropic background.
This is the dark matter component of the unresolved gamma-ray background (UGRB) --- also called the isotropic diffuse gamma-ray background (IGRB) --- whose total energy spectrum between 100 MeV and 820 GeV has been measured by the \textit{Fermi}-LAT~\cite{Fermi-LAT:2014ryh}. Well-established astrophysical source populations contribute their own guaranteed fraction to the UGRB: unresolved blazars (HSPs and LISPs) dominate at high energies and account for perhaps 20--30\% of the total intensity, while star-forming galaxies, misaligned AGNs, and millisecond pulsars add further contributions across the full energy range~\cite{Fornasa:2015qua}.
Because the exact composition of the UGRB is unknown, any dark matter signal must compete with and be separated from these astrophysical backgrounds.

The key challenge is spectral degeneracy.
At energies accessible to \textit{Fermi}-LAT ($0.1$--$300$ GeV), the spectra of heavy WIMPs and certain blazar populations overlap enough that energy information alone cannot distinguish them.
The annihilation flux is additionally sensitive to highly uncertain quantities: the concentration-mass relation of dark matter halos extrapolated to sub-resolution masses, the abundance of substructure, and the choice of minimum halo mass --- all of which enter through the so-called boost factor $B(M, z)$ that encodes how much the clumpy distribution of subhalos enhances the signal above a smooth halo assumption~\cite{Fornasa:2015qua}.
The boost factor itself can vary by orders of magnitude depending on these assumptions.
A purely spectral analysis of the UGRB therefore constrains combinations of particle physics parameters and astrophysical modeling choices that are difficult to disentangle.

Spatial anisotropies provide the orthogonal handle.
The UGRB is not perfectly isotropic: each contributing population imprints spatial fluctuations whose statistical properties --- the angular power spectrum $C_\ell$ --- are determined by the clustering properties of the sources and their redshift distribution.
Dark matter annihilation traces the squared density field of the cosmic web, which is more clustered on small angular scales and at lower redshifts than the blazars that constitute most of the astrophysical background~\cite{Ando:2005xg,Fermi-LAT:2012pez}. The \textit{Fermi}-LAT Collaboration provided the first measurement of the UGRB angular power spectrum, detecting an approximately scale-invariant (Poisson) component consistent with unresolved point sources, but without clear evidence for a dark matter contribution~\cite{Fermi-LAT:2012pez}.
The auto-correlation power spectrum, however, suffers from a fundamental limitation: the signal from dark matter is always mixed with the shot noise and correlated fluctuations of astrophysical sources that contribute at the same angular scales, and it lacks any redshift information.

Cross-correlating the gamma-ray map with an external gravitational tracer resolves both problems.
A galaxy catalog provides a redshift-selected map of large-scale structure that correlates with the dark matter distribution at specific cosmic distances.
The cross-power spectrum between the gamma-ray sky and the galaxy density field vanishes by construction for any source population that does not cluster with the selected galaxies --- this is the noise-cancellation advantage that makes cross-correlation a qualitatively more powerful probe than auto-correlation~\cite{Fornengo:2013rga,Camera:2012cj}.
The dark matter annihilation signal peaks at low redshift (because the flux is proportional to $\rho^2$ and thus weights dense, nearby structures more heavily), precisely where a catalog of local galaxies has the highest completeness.
By choosing a catalog with $z \lesssim 0.1$, one can arrange for the dark matter window function to overlap substantially with the galaxy window function while the blazar window function --- whose redshift distribution peaks at $z \sim 0.3$--0.4 --- mostly does not.
The angular cross-power spectrum developed formally in Section~\ref{sec:3.5} is the mathematical realization of this strategy.

Figure~\ref{fig:window_main} illustrates the core idea.
The dark matter window function $W_\gamma^\mathrm{DM}(z)$, which describes the redshift distribution of the gamma-ray emission, is sharply peaked at $z < 0.1$ due to the $\rho^2$ density weighting of the annihilation rate.
Local galaxy catalogs such as 2MASS, covering $z \lesssim 0.1$--0.2, have window functions that closely follow this peak.
In contrast, blazars at energies of 50 GeV are distributed over $z \sim 0.1$--0.4, and at 1 TeV they are further suppressed at high redshift by EBL absorption, but still do not overlap strongly with the very local galaxy distribution.
The cross-correlation signal between the gamma-ray sky and local galaxies is therefore dominated by whatever gamma-ray emission originates in the same volume as the galaxies --- precisely the regime where dark matter is the dominant cosmological source.
